% GENERATED FILE. Edit canonical lessons, then run npm run generate:books.
% canonical-source-hash: fnv1a64:7c8263adb45b85d3
% canonical-lessons: SA-C171-sandesah, SA-C171-suvidha, SA-C171-vyavastha, SA-C171-karyakramah, SA-C171-pratiyogita

\chapter{Message, Convenience, Arrangement, Programme, Competition}
\label{ch:sa-message-convenience-arrangement-programme-competition}
\hlchaptermodality{171}

\begin{chapteropening}
\textbf{By the end of this chapter:} \emph{I can say a message, a convenience, an arrangement, a programme and a competition.}

Complete the last lesson of chapter 171: I can say a message, a convenience, an arrangement, a programme and a competition.
\end{chapteropening}

\section[\texorpdfstring{\sk{सन्देशः} (sandeśaḥ)}{सन्देशः (sandeśaḥ)}]{\sk{सन्देशः} (sandeśaḥ) — a message}

\label{lesson:SA-C171-sandesah}



\begin{quote}
Before the new one: say the Sanskrit for a livelihood, then the Sanskrit for service.
\end{quote}

\subsection*{You'll want to know: \sk{सन्देशः}}
\textbf{\sk{सन्देशः}} — \emph{sandeśaḥ} — \textquotedblleft{}a message\textquotedblright{}.

\begin{grammarlens}[title={What you've built}]
One more word for this chapter.
\end{grammarlens}

\subsection*{Your turn}
\begin{itemize}
\raggedright
  \item \emph{Say it:} \emph{sandeśaḥ}
  \item \emph{Say it:} \emph{sandeśaḥ}, once more
  \item \emph{Say it:} say \emph{sandeśaḥ}
  \item \emph{Recall:} say the Sanskrit for a livelihood, then the Sanskrit for service, then say \emph{sandeśaḥ} again
\end{itemize}

\begin{tcolorbox}[breakable,colback=teal!4,colframe=teal!35!black,title={Before you move on}]
What does \sk{सन्देशः} mean? (\textquotedblleft{}A message\textquotedblright{}.) Say it once more.
\end{tcolorbox}
\section[\texorpdfstring{\sk{सुविधा} (suvidhā)}{सुविधा (suvidhā)}]{\sk{सुविधा} (suvidhā) — a convenience, a facility}

\label{lesson:SA-C171-suvidha}



\begin{quote}
Before the new one: say the Sanskrit for service, then the Sanskrit for a message.
\end{quote}

\subsection*{You'll want to know: \sk{सुविधा}}
\textbf{\sk{सुविधा}} — \emph{suvidhā} — \textquotedblleft{}a convenience, a facility\textquotedblright{}.

\begin{grammarlens}[title={What you've built}]
One more word for this chapter.
\end{grammarlens}

\subsection*{Your turn}
\begin{itemize}
\raggedright
  \item \emph{Say it:} \emph{suvidhā}
  \item \emph{Say it:} \emph{suvidhā}, once more
  \item \emph{Say it:} say \emph{suvidhā}
  \item \emph{Recall:} say the Sanskrit for service, then the Sanskrit for a message, then say \emph{suvidhā} again
\end{itemize}

\begin{tcolorbox}[breakable,colback=teal!4,colframe=teal!35!black,title={Before you move on}]
What does \sk{सुविधा} mean? (\textquotedblleft{}A convenience, a facility\textquotedblright{}.) Say it once more.
\end{tcolorbox}
\section[\texorpdfstring{\sk{व्यवस्था} (vyavasthā)}{व्यवस्था (vyavasthā)}]{\sk{व्यवस्था} (vyavasthā) — an arrangement}

\label{lesson:SA-C171-vyavastha}



\begin{quote}
Before the new one: say the Sanskrit for a message, then the Sanskrit for a convenience.
\end{quote}

\subsection*{You'll want to know: \sk{व्यवस्था}}
\textbf{\sk{व्यवस्था}} — \emph{vyavasthā} — \textquotedblleft{}an arrangement\textquotedblright{}.

\begin{grammarlens}[title={What you've built}]
One more word for this chapter.
\end{grammarlens}

\subsection*{Your turn}
\begin{itemize}
\raggedright
  \item \emph{Say it:} \emph{vyavasthā}
  \item \emph{Say it:} \emph{vyavasthā}, once more
  \item \emph{Say it:} say \emph{vyavasthā}
  \item \emph{Recall:} say the Sanskrit for a message, then the Sanskrit for a convenience, then say \emph{vyavasthā} again
\end{itemize}

\begin{tcolorbox}[breakable,colback=teal!4,colframe=teal!35!black,title={Before you move on}]
What does \sk{व्यवस्था} mean? (\textquotedblleft{}An arrangement\textquotedblright{}.) Say it once more.
\end{tcolorbox}
\section[\texorpdfstring{\sk{कार्यक्रमः} (kāryakramaḥ)}{कार्यक्रमः (kāryakramaḥ)}]{\sk{कार्यक्रमः} (kāryakramaḥ) — a programme, an event}

\label{lesson:SA-C171-karyakramah}



\begin{quote}
Before the new one: say the Sanskrit for a convenience, then the Sanskrit for an arrangement.
\end{quote}

\subsection*{You'll want to know: \sk{कार्यक्रमः}}
\textbf{\sk{कार्यक्रमः}} — \emph{kāryakramaḥ} — \textquotedblleft{}a programme, an event\textquotedblright{}.

\begin{grammarlens}[title={What you've built}]
One more word for this chapter.
\end{grammarlens}

\subsection*{Your turn}
\begin{itemize}
\raggedright
  \item \emph{Say it:} \emph{kāryakramaḥ}
  \item \emph{Say it:} \emph{kāryakramaḥ}, once more
  \item \emph{Say it:} say \emph{kāryakramaḥ}
  \item \emph{Recall:} say the Sanskrit for a convenience, then the Sanskrit for an arrangement, then say \emph{kāryakramaḥ} again
\end{itemize}

\begin{tcolorbox}[breakable,colback=teal!4,colframe=teal!35!black,title={Before you move on}]
What does \sk{कार्यक्रमः} mean? (\textquotedblleft{}A programme, an event\textquotedblright{}.) Say it once more.
\end{tcolorbox}
\section[\texorpdfstring{\sk{प्रतियोगिता} (pratiyogitā)}{प्रतियोगिता (pratiyogitā)}]{\sk{प्रतियोगिता} (pratiyogitā) — a competition}

\label{lesson:SA-C171-pratiyogita}



\begin{quote}
Before the new one: say the Sanskrit for an arrangement, then the Sanskrit for a programme.
\end{quote}

\subsection*{You'll want to know: \sk{प्रतियोगिता}}
\textbf{\sk{प्रतियोगिता}} — \emph{pratiyogitā} — \textquotedblleft{}a competition\textquotedblright{}.

\begin{grammarlens}[title={What you've built}]
One more word for this chapter.
\end{grammarlens}

\subsection*{Your turn}
\begin{itemize}
\raggedright
  \item \emph{Say it:} \emph{pratiyogitā}
  \item \emph{Say it:} \emph{pratiyogitā}, once more
  \item \emph{Say it:} say \emph{pratiyogitā}
  \item \emph{Recall:} say the Sanskrit for an arrangement, then the Sanskrit for a programme, then say \emph{pratiyogitā} again
\end{itemize}

\begin{tcolorbox}[breakable,colback=teal!4,colframe=teal!35!black,title={Before you move on}]
What does \sk{प्रतियोगिता} mean? (\textquotedblleft{}A competition\textquotedblright{}.) Say it once more.
\end{tcolorbox}
